% Kamp III.2, Lemma 4: conjunction by finite fusion of interval patterns.

\subsubsection{Conjunction of elementary interval conditions}
\label{sec:p09-kamp-iii2-lemma4}

Kamp's next step combines two finite interval descriptions with the same
ordered endpoints. Their internal witnesses can interleave or coincide.
Lemma~4 accounts for every such arrangement by a finite disjunction of
elementary interval formulas. The source introduces the classes and the
property on printed pp.~67--68; the complete proof runs from p.~73 to the
explicit end of Lemma~4 on p.~78 \cite[ch.~III, pp.~67--68, 73--78]{Kamp1968}.

\begin{hbdefinition}{def:p09-kamp-elementary-intervals}{Elementary past and future interval classes}
Use the basic classes of \cref{def:p09-kamp-basic-classes}, the class
\(\operatorname{E}\) of \cref{def:p09-kamp-elementary-class}, and proper
unary-predicate substitution from \cref{def:p09-kamp-proper-substitution}.
For distinct variable names \(x,y\), a member of
\(\operatorname{Ep}(x,y)\) is obtained from a member of
\(\operatorname{Bp}(x,y)\) by a finite sequence of proper substitutions
whose replacements belong to \(\operatorname{E}\). The corresponding
construction from \(\operatorname{Bf}(x,y)\) defines
\(\operatorname{Ef}(x,y)\).

More literally, the source writes the successive targets as
\(P_0,P_1,\ldots,P_N\), for some finite \(N\), with a replacement in
\(\operatorname{E}\) for each target. Identity replacements
\(P_i(w)\) allow unused positions in such a list. The equivalence of this
initial-segment description and arbitrary finite substitution histories is
proved in the annex. Replacements are applied successively in the stated
order; a replacement inserted at one step can therefore be affected by a
later step. Within a single step, its newly inserted atoms are not visited
again. Repeated and shared predicate letters are allowed.

As in the preceding substitution construction, bound variables are chosen
fresh before inserting a replacement, and formulas obtained by different
such choices are compared up to alphabetic variation. Every elementary
interval formula has exactly \(x,y\) free. Neither binary class inserts
an outer assertion \(x<y\).
These are Kamp's elementary past and elementary future classes, distinct
from the one-variable class \(\operatorname{E}\)
\cite[ch.~III, p.~67]{Kamp1968}.
\end{hbdefinition}

\begin{hbdefinition}{def:p09-kamp-labelled-intervals}{Elementary labels on a finite interval word}
Fix a variable \(w\). An \emph{elementary labelled word} is a finite list
\[
 W=(\alpha_0,\beta_1,\alpha_1,\ldots,\beta_m,\alpha_m),
 \qquad m\geq1,
\]
whose entries belong to \(\operatorname{E}(w)\). Repetitions are
permitted and retain their separate positions. For a label \(\theta\)
and a variable \(v\), write \(\theta(v)\) for the capture-avoiding
replacement \(\theta[v/w]\), after making its bound variables fresh.
This is a formula abbreviation, not application of a new predicate symbol.

Let \(\mathcal B_p(W;x,y)\) be the basic past pattern with these
one-variable formulas in its successive point and interval positions. Its
explicit recursion is
\[
\begin{split}
 \mathcal B_p((\alpha_0,\beta_1,\alpha_1);x,y)
 &\defeq \alpha_0(x)\land
   \forall v\,(x<v<y\to\beta_1(v))\land\alpha_1(y),\\
 \mathcal B_p(W;x,y)
 &\defeq\alpha_0(x)\land\exists z\bigl(x<z<y\land
   \forall v\,(x<v<z\to\beta_1(v))\\
 &\hspace{43mm}{}\land\mathcal B_p(W^-;z,y)\bigr)\quad(m>1),
\end{split}
\]
where \(W^-=(\alpha_1,\beta_2,\alpha_2,\ldots,\beta_m,\alpha_m)\).
The notation \(a<b<c\) abbreviates \(a<b\land b<c\).
Every displayed bound variable is fresh from the endpoints and from all
labels at that occurrence. The future version has the same one-interval
base and the recursion
\[
\begin{split}
 \mathcal B_f(W;x,y)\defeq
 \exists z\bigl(x<z<y\land\mathcal B_f(W^+;x,z)
 &\land\forall v\,(z<v<y\to\beta_m(v))\bigr)\\
 &{}\land\alpha_m(y)\quad(m>1),
\end{split}
\]
where \(W^+=(\alpha_0,\beta_1,\ldots,\beta_{m-1},\alpha_{m-1})\).

In a linear structure \(\mathcal M\), a realization of \(W\) is a
strictly increasing sequence \(a_0<\cdots<a_m\) such that
\(\alpha_i\) holds with \(w=a_i\), and \(\beta_i\) holds with
\(w=t\) for every \(a_{i-1}<t<a_i\). A label is evaluated in the
whole structure with its stated free argument. Its internal quantifiers
are not restricted to the displayed interval. Only the outer interval
quantifier in the pattern has that restriction.
\end{hbdefinition}

\begin{hblemma}{thm:p09-kamp-elementary-intervals}{proved=annex:p9.kamp.elementary.intervals}{Elementary words, substitution and the interval remarks}
Every elementary interval formula has an elementary labelled word obtained
by following its finite substitution history at each original predicate
occurrence. Conversely, every \(\mathcal B_p(W;x,y)\) has a
\(\operatorname{Ep}(x,y)\) presentation, and every
\(\mathcal B_f(W;x,y)\) has an \(\operatorname{Ef}(x,y)\)
presentation, with only fresh bound names possibly differing.
Whenever \(a<b\), either presentation holds at \((a,b)\) exactly when
\(W\) has a realization beginning at \(a\) and ending at \(b\).
These assertions hold over \(\mathsf{Lin}\).

In particular the source's remarks R5--R6 hold: if
\(\phi\in\operatorname{Ep}(x,y)\), its guarded past closure
\(\exists x\,(x<y\land\phi)\) has an \(\operatorname{E}(y)\)
presentation; if \(\phi\in\operatorname{Ef}(x,y)\), its guarded
future closure \(\exists y\,(x<y\land\phi)\) has an
\(\operatorname{E}(x)\) presentation. Proper substitution of another
member of \(\operatorname{E}\) preserves either elementary interval
class. These are statements about the displayed constructions with allowed
bound renaming, and do not add semantic-equivalence closure to the grammar
of \(\operatorname{E}\) \cite[ch.~III, pp.~67--68]{Kamp1968}.
\end{hblemma}

\begin{hbdefinition}{def:p09-kamp-interval-property}{Kamp's interval property \(I(x,y)\)}
Let \(x,y\) be distinct variable names. An \(L\)-formula \(\phi\)
with no free variables other than possibly
\(x,y\) has \emph{property \(I(x,y)\)} if there is a finite
disjunction of formulas \(\theta_i\in\operatorname{Ep}(x,y)\) such that
\[
 \mathsf{Com}\models
 x<y\longrightarrow
             \left(\phi\longleftrightarrow\bigvee_i\theta_i\right).
\]
The equivalence concerns every unary-predicate interpretation and every
assignment on every complete linear order. It makes no assertion about
assignments with \(x\not<y\). No condition that both variables actually
occur freely in \(\phi\) is imposed.

A disjunction may be taken nonempty: a false interval formula is represented
by \(\mathcal B_p((F^K_w,T^K_w,T^K_w);x,y)\), with the source
truth and falsehood formulas of \cref{def:p09-kamp-elementary-class}.
Thus an empty disjunction, if used as a finite Boolean abbreviation, can
be replaced by this one disjunct. This convention introduces no new
primitive connective \cite[ch.~III, p.~68, display~(13)]{Kamp1968}.
\end{hbdefinition}

\begin{hblemma}{thm:p09-kamp-interval-substitution}{proved=annex:p9.kamp.interval.substitution}{Substitution and one-point instances of the interval property}
If \(\phi\) has property \(I(x,y)\) and
\(\psi\in\operatorname{E}\), then
\(\operatorname{psub}_{P_i}(\phi,\psi)\) has property \(I(x,y)\).
If \(\theta\in\operatorname{E}(x)\) and \(y\) is a different
variable name, then \(\theta\) has both property \(I(x,y)\) and
property \(I(y,x)\). These are the source's remarks R7--R8; the annex
proves their substitution semantics and both endpoint placements
\cite[ch.~III, p.~68]{Kamp1968}.
\end{hblemma}

\begin{hbdefinition}{def:p09-kamp-fusions}{Fusions of two endpoint sequences}
Let the two input words be
\[
 W=(\alpha_0,\beta_1,\ldots,\beta_m,\alpha_m),\qquad
 V=(\gamma_0,\delta_1,\ldots,\delta_n,\gamma_n),
 \qquad m,n\geq1.
\]
Give their point positions distinct tags
\((e_0,\ldots,e_m)\) and \((f_0,\ldots,f_n)\). Tags record
occurrences, even when the formulas attached to two positions coincide.

A \emph{total preorder} on these tags is a reflexive, transitive binary
relation \(F\) comparing every two tags. Put
\(a\sim_F b\) when both \(aFb\) and \(bFa\) hold. A
\emph{fusion} is such a relation whose restriction to each input sequence
is exactly that sequence's reflexive linear order, and which satisfies
\[
                      e_0\sim_F f_0,
                \qquad e_m\sim_F f_n.
\]
Thus different positions of the same input cannot coincide, while a
position of one input may coincide with a position of the other. The
endpoint identifications are mandatory. Write \(\mathcal F_{m,n}\)
for the finite set of these fusions.

For a fusion \(F\), enumerate its equivalence classes in their strict
order as \(C_1,\ldots,C_{r_F}\). Define \(h_F(a)=j\) when
\(a\in C_j\). In particular,
\[
 h_F(e_0)=h_F(f_0)=1,\qquad
 h_F(e_m)=h_F(f_n)=r_F,\qquad r_F\geq2.
\]
Kamp calls a total preorder a \emph{quasi-ordering}, a finite sequence
without repetition a \emph{pure sequence}, and two such sequences with
disjoint members a \emph{pure pair}. His general fusion preserves the
two input orders; the endpoint identifications select the fusions used
in Lemma~4. His \emph{independent} members of the second
sequence are those not equivalent to any member of the first. Taking all
members of the first sequence and those independent members of the second
selects exactly one representative of each class. Hence his fusion length
is the number \(r_F\) of classes used here
\cite[ch.~III, pp.~74--75]{Kamp1968}.
\end{hbdefinition}

\begin{hbmethod}{def:p09-kamp-fusion-construction}{The formula associated with a fusion}
In addition to the point tags, give the interval positions tags
\(b_1,\ldots,b_m,d_1,\ldots,d_n\), all distinct. Let \(\mathcal T\)
be the set of all point and interval tags. Attach the labels
\[
 \theta_{e_i}=\alpha_i,\quad \theta_{b_i}=\beta_i,\quad
 \theta_{f_j}=\gamma_j,\quad \theta_{d_j}=\delta_j.
\]
For a fixed fusion \(F\), there are \(2r_F-1\) alternating slots:
slot \(2k-1\) is its \(k\)-th point, and slot \(2k\) is the open
gap between its \(k\)-th and \((k+1)\)-st points. Define the set of
slots required by each input tag:
\[
\begin{aligned}
 f_F(e_i)&\defeq\{2h_F(e_i)-1\},&
 f_F(f_j)&\defeq\{2h_F(f_j)-1\},\\
 f_F(b_i)&\defeq
   \{s:2h_F(e_{i-1})-1<s<2h_F(e_i)-1\},\\
 f_F(d_j)&\defeq
   \{s:2h_F(f_{j-1})-1<s<2h_F(f_j)-1\}.
\end{aligned}
\]
The variable \(s\) ranges over integers in \(\{1,\ldots,2r_F-1\}\).
The notation \(f_F\) is a set-valued map on tags; the tag \(f_j\)
is a point position of the second input. These have different declared
roles. In particular an interval tag can require both even gap slots and
odd point slots strictly between its endpoints.

For each slot put
\[
 \rho_{F,s}(w)\defeq
       \bigwedge_{\substack{\tau\in\mathcal T\\s\in f_F(\tau)}}\theta_\tau(w).
\]
Use the tag order
\(e_0,b_1,e_1,\ldots,b_m,e_m,f_0,d_1,f_1,\ldots,d_n,f_n\)
to fix the order of this finite conjunction.
Every slot has at least one required tag, as proved in the annex, so its
label belongs to \(\operatorname{E}(w)\). All listed conjuncts are
required; no negative requirement is imposed on any unlisted label.
Finally set
\[
 \chi_F(x,y)\defeq
  \mathcal B_p((\rho_{F,1},\ldots,\rho_{F,2r_F-1});x,y).
\]
It has an \(\operatorname{Ep}(x,y)\) presentation by
\cref{thm:p09-kamp-elementary-intervals}. Enumerate all fusions and
return \(\bigvee_{F\in\mathcal F_{m,n}}\chi_F(x,y)\).
This is the slot construction of source pp.~75--76, with point positions
and alternating slots indexed separately.
\end{hbmethod}

\begin{hblemma}{thm:p09-kamp-iii2-lemma4}{proved=annex:p9.kamp.iii2.lemma4}{Kamp's conjunction lemma for \(I(x,y)\)}
If \(\phi,\psi\) have property \(I(x,y)\), then \(\phi\land\psi\)
has property \(I(x,y)\).

The elementary construction used in the proof has the stronger, explicit
linear-order identity
\[
\begin{split}
 \mathsf{Lin}\models x<y\longrightarrow
 \Bigl(&\mathcal B_p(W;x,y)\land\mathcal B_p(V;x,y)\\
       &\longleftrightarrow
                   \bigvee_{F\in\mathcal F_{m,n}}\chi_F(x,y)\Bigr).
\end{split}
\]
For given finite \(I(x,y)\) representations, distribute their conjunction
over every pair of disjuncts, construct their labelled words, and apply this
fusion method to each pair. Restricting the displayed identity to
\(\mathsf{Com}\) gives the required representation of
\(\phi\land\psi\). All constructions are finite and use only
predicate letters in the input elementary labels; temporary fresh letters
used to exhibit membership are eliminated
\cite[ch.~III, lemma~4, pp.~73--78]{Kamp1968}.
\end{hblemma}

\begin{hbexample}{def:p09-kamp-fusion-three-cases}{Two interior points: either order or coincidence}
Take the two words
\(W=(A_0,B_1,A_1,B_2,A_2)\) and
\(V=(C_0,D_1,C_1,D_2,C_2)\), with all entries in
\(\operatorname{E}(w)\). Write \(a\) and \(c\) for their respective
interior witnesses. The three fusions have the following labelled words;
conjunction here combines labels at the same parameter \(w\):
\[
\begin{array}{c|l}
 a<c & (A_0\land C_0,\ B_1\land D_1,\ A_1\land D_1,\
           B_2\land D_1,\\[-2pt]
     & \hspace{12mm} C_1\land B_2,\ B_2\land D_2,\ A_2\land C_2)\\[3pt]
 a=c & (A_0\land C_0,\ B_1\land D_1,\ A_1\land C_1,\
           B_2\land D_2,\ A_2\land C_2)\\[3pt]
 c<a & (A_0\land C_0,\ B_1\land D_1,\ C_1\land B_1,\
           B_1\land D_2,\\[-2pt]
     & \hspace{12mm} A_1\land D_2,\ B_2\land D_2,\ A_2\land C_2).
\end{array}
\]
The coincident case cannot be discarded: on the three-point order
\(0<1<2\), take both middle labels true only at 1 and every other
label true everywhere. The two input patterns hold at \((0,2)\), and
only the coincident fusion can be realized.

An interval label must also hold at an inserted interior point. On
\(0<1<2<3\), take \(A_1\) true only at 1, \(C_1\) true only at 2,
\(B_2\) false everywhere, and all other labels true everywhere. The
first input fails because its interval \((1,3)\) contains 2. In the
\(a<c\) word the conjunct \(B_2\) at the \(C_1\) position detects
this failure; the two neighboring open gaps are empty and cannot detect it.
The annex verifies these cases and the one-interval and endpoint cases.
\end{hbexample}

\paragraph{Order scope and the next source result.}
Property \(I(x,y)\) is defined with \(\mathsf{Com}\) on p.~68, and
the convention immediately before Lemma~4 assigns that same scope to
unqualified equivalence. The opening proof on p.~73 nevertheless prints
\(\mathsf{Lin}\) on the representation premises. The argument here
uses those premises with their defining \(\mathsf{Com}\) scope and
proves the elementary fusion identity separately over \(\mathsf{Lin}\).
Its restriction to complete orders is sufficient; no converse transfer
from a complete-order equivalence to all linear orders is assumed.

The next result is Lemma~5, beginning on printed p.~78. It concerns the
guarded existential composition through an interior point \(z\), with
\(\operatorname{E}(z)\) conditions at that point and interval properties
\(I(x,z)\), \(I(z,y)\) on either side. That construction and the later
lemmas remain separate from the conjunction result proved here
\cite[ch.~III, pp.~68, 73, 78--79]{Kamp1968}.
